La comunicación entre módulos es la definida en el diagrama de flujo que se presenta en la Figura 7.1.

El núcleo principal de la configuración entre módulos son los archivos \textit{main.py} y \textit{config.py}. En el archivo \textit{main} se decide
la acción que se va a ejecutar y se llama a una función que lleva a cabo la tarea requerida. Previamente se cambian valores del archivo config, como por ejemplo, la función de pérdida.

\begin{itemize}
    \item 1. CLI / GUI modifica configuración
    \item 2. Se cargan datos o audio externo.
    \item 3. Se calcula STFT.
    \item 4. El modelo estima máscaras.
    \item 5. \textit{run\_inference} reconstruye audio.
    \item 6. \textit{deployment} guarda stems.
    \item 7. Evaluación calcula métricas.
    \item 8. Resultados se guardan en carpeta de checkpoints/JSON/gráficas/tablas.

\end{itemize}
